\documentclass[UTF8]{ctexart}
\usepackage{xeCJK}
\usepackage{geometry}
\geometry{left=2cm,right=2cm,top=2.5cm,bottom=2.5cm}
\setlength{\headheight}{12.7pt}

\usepackage{titlesec}
\titleformat{\section}
  {\normalfont\large\bfseries}
  {\thesection}
  {1em}
  {}
\titlespacing*{\section}
  {0pt}
  {3.5ex plus 1ex minus .2ex}
  {2.3ex plus .2ex}

\usepackage{amsmath}
\usepackage{amssymb}
\usepackage{amsthm}
\usepackage{enumitem}
\setlist[enumerate]{itemsep=0pt,topsep=0pt,parsep=0pt,leftmargin=0pt,itemindent=2.5em}
\setlist[itemize]{itemsep=0pt,topsep=0pt,parsep=0pt,leftmargin=0pt,itemindent=2.5em}
\usepackage{fancyhdr}
\pagestyle{fancy}
\fancyhf{}
\usepackage{graphicx}
\usepackage{hyperref}
\usepackage{indentfirst}
\usepackage{lmodern}


\begin{document}
\setlength{\abovedisplayskip}{1pt}
\setlength{\belowdisplayskip}{1pt}

\section{一阶语言}

\hypertarget{sec:定义1.1}{}
\noindent\textbf{定义1.1 (一阶语言)}

称$\mathcal{L}=(\mathcal{V},\mathcal{C},
(n_f)_{f\in\mathcal{F}},(n_r)_{r\in\mathcal{R}})$是一个\textbf{一阶语言}, 如果:
\begin{enumerate}
    \item $\mathcal{V}$是无限集,
    \item $\mathcal{V},\mathcal{C},\mathcal{F},\mathcal{R}$两两不相交,
    \item 对任意的$f\in\mathcal{F}$有$n_f\in\mathbb{N}^+$,
    对任意的$r\in\mathcal{R}$有$n_r\in\mathbb{N}^+$.
\end{enumerate}

\vspace{10pt}

分别称$\mathcal{V},\mathcal{C},\mathcal{F},\mathcal{R}$为
变元符号集、常元符号集、函数符号集和关系符号集.

称$n_f$为$f$的元数, 称$n_r$为$r$的元数.

\vspace{10pt}

此时记$|\mathcal{L}|\overset{\mathrm{def}}{=}
|\mathcal{V}\cup\mathcal{C}\cup\mathcal{F}\cup\mathcal{R}|$,
称之为$\mathcal{L}$的基数. 显然这一定是无限基数.

\vspace{10pt}

\hypertarget{sec:定义1.2}{}
\noindent\textbf{定义1.2 (一阶语言的字母表)}

给定一阶语言$\mathcal{L}=(\mathcal{V},\mathcal{C},
(n_f)_{f\in\mathcal{F}},(n_r)_{r\in\mathcal{R}})$, 称
\[
    \Sigma_{\mathcal{L}}\overset{\mathrm{def}}{=}
    \mathcal{V}\cup\mathcal{C}\cup\mathcal{F}\cup\mathcal{R}\cup\mathrm{Log}
\]
为$\mathcal{L}$的\textbf{字母表}, 其中
$\mathrm{Log}=\{\ \neg\ ,\ \to\ ,\ \exists\ ,\ \equiv\ ,\ (\ ,\ )\ \}$.

默认$\mathrm{Log}$也与$\mathcal{L}$中的各集合均不相交.

\vspace{10pt}

记$\Sigma_{\mathcal{L}}^*$为$\Sigma_{\mathcal{L}}$中符号组成的所有字符串的集合.

\vspace{10pt}

\hypertarget{sec:定理1.1}{}
\noindent\textbf{定理1.1}

对任意的一阶语言$\mathcal{L}$, $|\Sigma_{\mathcal{L}}^*|=|\mathcal{L}|$.

\noindent{\kaishu 证明.}

显然$|\Sigma_{\mathcal{L}}|=|\mathcal{L}|$.
对任意的自然数$n$, 把$\Sigma_{\mathcal{L}}$中长度为$n$的字符串的集合记为
$\Sigma_{\mathcal{L}}^n$, 则有$\displaystyle\Sigma_{\mathcal{L}}^*=
\bigcup_{n\in\mathbb{N}}\Sigma_{\mathcal{L}}^n$.
\begin{enumerate}
    \item 显然$|\mathbb{N}|=\omega\leqslant|\mathcal{L}|$,
    \item 对任意的自然数$n$, 显然$|\Sigma_{\mathcal{L}}^n|=
    |\Sigma_{\mathcal{L}}|^n\leqslant|\Sigma_{\mathcal{L}}|=|\mathcal{L}|$,
\end{enumerate}
\vspace{3pt}
于是$\displaystyle|\Sigma_{\mathcal{L}}^*|=
\left|\bigcup_{n\in\mathbb{N}}\Sigma_{\mathcal{L}}^n\right|\leqslant|\mathcal{L}|$.
同时显然$|\Sigma_{\mathcal{L}}^*|\geqslant|\mathcal{L}|$, 所以二者相等.
\hfill $\qedsymbol$

\vspace{10pt}

\hypertarget{sec:定义1.3}{}
\noindent\textbf{定义1.3 (子语言)}

给定一阶语言$\mathcal{L}_1,\mathcal{L}_2$. 如果$\mathcal{L}_1$内的所有符号
也都是$\mathcal{L}_2$的对应符号, 并且函数和关系符号的元数相同, 就称$\mathcal{L}_1$
是$\mathcal{L}_2$的\textbf{子语言}, 记作$\mathcal{L}_1\subset\mathcal{L}_2$.





\newpage

\section*{附录1\ \ 归纳定义与归纳证明}
\addcontentsline{toc}{section}{(附录) 归纳定义与归纳证明}

\hypertarget{sec:定义A1.1}{}
\noindent\textbf{定义A1.1 (归纳定义)}

给定集合$U$以及:
\begin{enumerate}
    \item 初始集$B\subset U$,
    \item 归纳规则: 一个索引集$I$, 并且每个$i\in I$对应一个正整数$n_i$和一个映射
    \[
        f_i:U^{n_i}\to U.
    \]
\end{enumerate}
那么可以对每个自然数$k$递归定义$G_k$如下:
\begin{enumerate}
    \item $G_0=B$,
    \item 对每个自然数$k$归纳定义
    \[
        G_{k+1}=\left(\bigcup_{i\in I}\{f_i(u_1,\cdots,u_{n_i})
        \mid u_1,\cdots,u_{n_i}\in G_k\}\right)\cup G_k,
    \]
\end{enumerate}
最后记$\displaystyle G=\bigcup_{k\in\mathbb{N}}G_k$,
称之为由初始集和归纳规则归纳定义的目标集. 称$G_k$为其第$k$层.

\vspace{10pt}

依定义, 任意选定$i\in I$和$u_1,\cdots,u_{n_i}\in G$,
显然有$f_i(u_1,\cdots,u_{n_i})\in G$.

\vspace{10pt}

\hypertarget{sec:定理A1.1}{}
\noindent\textbf{定理A1.1}

给定集合$U$, 初始集$B$和归纳规则$(n_i)_{i\in I},(f_i)_{i\in I}$,
记其归纳定义的目标集为$G$, 有
\[
    |G|\leqslant\max\{|B|,|I|,\omega\}.
\]

\noindent{\kaishu 证明.}

记$\kappa=\max\{|B|,|I|,\omega\}$, 归纳证明每一层$|G_k|\leqslant\kappa$.
首先$|G_0|=|B|\leqslant\kappa$; 其次假设$|G_k|\leqslant\kappa$, 则
\begin{enumerate}
    \item $|I|\leqslant\kappa$,
    \item 对每个$i\in I$, $|\{f_i(u_1,\cdots,u_{n_i})\mid
    u_1,\cdots,u_{n_i}\in G_k\}|\leqslant|G_k|^{n_i}\leqslant\kappa^{n_i}=\kappa$,
\end{enumerate}
所以$\displaystyle\left|\bigcup_{i\in I}\{f_i(u_1,\cdots,u_{n_i})\mid
u_1,\cdots,u_{n_i}\in G_k\}\right|\leqslant\kappa$,
进一步$|G_{k+1}|\leqslant\kappa$. 归纳证明成立.

最后, 因为$|\mathbb{N}|\leqslant\kappa$并且对每个$k\in\mathbb{N}$有
$|G_k|\leqslant\kappa$, 所以$\displaystyle
|G|=\left|\bigcup_{k\in\mathbb{N}}G_k\right|\leqslant\kappa$. \hfill $\qedsymbol$

\vspace{10pt}

\hypertarget{sec:定理A1.2}{}
\noindent\textbf{定理A1.2 (归纳证明)}

给定集合$U$, 初始集$B$和归纳规则$(n_i)_{i\in I},(f_i)_{i\in I}$,
记其归纳定义的目标集为$G$.

给定命题$p(u)$. 如果
\begin{enumerate}
    \item $\forall u\in B:p(u)$,
    \item 对任意的$i\in I$有
    \[
        \forall u_1,\cdots,u_{n_i}\in G:p(u_1)\land\cdots\land p(u_{n_i})\to
        p(f_i(u_1,\cdots,u_{n_i}));
    \]
\end{enumerate}
那么$\forall u\in G:p(u)$.

\noindent{\kaishu 证明.}

只需归纳证明每一层的元素都满足命题$p$.

首先, 显然$G_0$的元素都满足命题$p$;
其次, 假设$G_k$的元素都满足命题$p$, 那么任取$u\in G_{k+1}$:
\begin{enumerate}
    \item 如果存在$i\in I$和$u_1,\cdots,u_{n_i}\in G_k$使得
    $u=f_i(u_1,\cdots,u_{n_i})$, 那么依假设可得$p(u)$,
    \item 如果$u\in G_k$, 依假设即得$p(u)$,
\end{enumerate}
所以$G_{k+1}$的元素也都满足命题$p$. {\hfill $\qedsymbol$}

\end{document}